% GENERATED FILE. Edit canonical lessons, then run npm run generate:books.
% canonical-source-hash: fnv1a64:53d5efc2981956fb
% canonical-lessons: JA-C494-tamanegi, JA-C494-kyuuri, JA-C494-kinoko, JA-C494-takenoko, JA-C494-shiitake

\chapter{Onions and Mushrooms}
\label{ch:ja-onions-and-mushrooms}
\hlchaptermodality{494}

\begin{chapteropening}
\textbf{By the end of this chapter:} \emph{I can name more vegetables.}

Complete the last lesson of chapter 494: I can name more vegetables.
\end{chapteropening}

\section[\texorpdfstring{\ja{たまねぎ} (tamanegi)}{たまねぎ (tamanegi)}]{\ja{たまねぎ} (tamanegi) — an onion}

\label{lesson:JA-C494-tamanegi}



\begin{quote}
Before the new one: say the Japanese for the washing-up, then the Japanese for to freeze (food).
\end{quote}

\subsection*{You'll want to know: \ja{たまねぎ}}
\textbf{\ja{たまねぎ}} — \emph{tamanegi} — \textquotedblleft{}an onion\textquotedblright{}.

\begin{quote}
\textbf{\ja{たまねぎを} \ja{きざんだ} \ja{とき、なきました。}} — \emph{tamanegi o kizanda toki, nakimashita} — \textquotedblleft{}When I chopped the onion, I cried.\textquotedblright{}
\end{quote}

\begin{cousinweb}
\textbf{\ja{たまねぎ}} (\emph{tamanegi}) is a native Japanese compound (wago), inherited rather than borrowed from Chinese: built from \emph{tama}, \textquotedblleft{}ball\textquotedblright{} and \emph{negi}, \textquotedblleft{}leek\textquotedblright{} (JMdict, KANJIDIC2).
\end{cousinweb}

\begin{grammarlens}[title={What you've built}]
One more word for describing what you lived and felt.
\end{grammarlens}

\subsection*{Your turn}
\begin{itemize}
\raggedright
  \item \emph{Say it:} \emph{tamanegi}
  \item \emph{Say it:} \emph{tamanegi}, once more
  \item \emph{Say it:} \emph{tamanegi o kizanda toki, nakimashita} — When I chopped the onion, I cried.
  \item \emph{Recall:} say the Japanese for the washing-up, then the Japanese for to freeze (food), then say \emph{tamanegi} again
\end{itemize}

\begin{tcolorbox}[breakable,colback=teal!4,colframe=teal!35!black,title={Before you move on}]
What does \ja{たまねぎ} mean? (\textquotedblleft{}An onion\textquotedblright{}.) Say it once more.
\end{tcolorbox}
\section[\texorpdfstring{\ja{きゅうり} (kyūri)}{きゅうり (kyūri)}]{\ja{きゅうり} (kyūri) — a cucumber}

\label{lesson:JA-C494-kyuuri}



\begin{quote}
Before the new one: say the Japanese for to freeze (food), then the Japanese for an onion.
\end{quote}

\subsection*{You'll want to know: \ja{きゅうり}}
\textbf{\ja{きゅうり}} — \emph{kyūri} — \textquotedblleft{}a cucumber\textquotedblright{}.

\begin{quote}
\textbf{\ja{なつは} \ja{きゅうりが} \ja{やすいです。}} — \emph{natsu wa kyūri ga yasui desu} — \textquotedblleft{}Cucumbers are cheap in summer.\textquotedblright{}
\end{quote}

\begin{cousinweb}
\textbf{\ja{きゅうり}} (\emph{kyūri}) is a native Japanese word (wago), inherited rather than borrowed from Chinese, and usually written in kana; its older origin is not settled (JMdict).
\end{cousinweb}

\begin{grammarlens}[title={What you've built}]
One more word for describing what you lived and felt.
\end{grammarlens}

\subsection*{Your turn}
\begin{itemize}
\raggedright
  \item \emph{Say it:} \emph{kyūri}
  \item \emph{Say it:} \emph{kyūri}, once more
  \item \emph{Say it:} \emph{natsu wa kyūri ga yasui desu} — Cucumbers are cheap in summer.
  \item \emph{Recall:} say the Japanese for to freeze (food), then the Japanese for an onion, then say \emph{kyūri} again
\end{itemize}

\begin{tcolorbox}[breakable,colback=teal!4,colframe=teal!35!black,title={Before you move on}]
What does \ja{きゅうり} mean? (\textquotedblleft{}A cucumber\textquotedblright{}.) Say it once more.
\end{tcolorbox}
\section[\texorpdfstring{\ja{きのこ} (kinoko)}{きのこ (kinoko)}]{\ja{きのこ} (kinoko) — a mushroom}

\label{lesson:JA-C494-kinoko}



\begin{quote}
Before the new one: say the Japanese for an onion, then the Japanese for a cucumber.
\end{quote}

\subsection*{You'll want to know: \ja{きのこ}}
\textbf{\ja{きのこ}} — \emph{kinoko} — \textquotedblleft{}a mushroom\textquotedblright{}.

\begin{quote}
\textbf{\ja{あきの} \ja{やまで} \ja{きのこを} \ja{みつけました。}} — \emph{aki no yama de kinoko o mitsukemashita} — \textquotedblleft{}I found mushrooms on the mountain in autumn.\textquotedblright{}
\end{quote}

\begin{cousinweb}
\textbf{\ja{きのこ}} (\emph{kinoko}) is a native Japanese word (wago), inherited rather than borrowed from Chinese, and usually written in kana; its older origin is not settled (JMdict).
\end{cousinweb}

\begin{grammarlens}[title={What you've built}]
One more word for describing what you lived and felt.
\end{grammarlens}

\subsection*{Your turn}
\begin{itemize}
\raggedright
  \item \emph{Say it:} \emph{kinoko}
  \item \emph{Say it:} \emph{kinoko}, once more
  \item \emph{Say it:} \emph{aki no yama de kinoko o mitsukemashita} — I found mushrooms on the mountain in autumn.
  \item \emph{Recall:} say the Japanese for an onion, then the Japanese for a cucumber, then say \emph{kinoko} again
\end{itemize}

\begin{tcolorbox}[breakable,colback=teal!4,colframe=teal!35!black,title={Before you move on}]
What does \ja{きのこ} mean? (\textquotedblleft{}A mushroom\textquotedblright{}.) Say it once more.
\end{tcolorbox}
\section[\texorpdfstring{\ja{たけのこ} (takenoko)}{たけのこ (takenoko)}]{\ja{たけのこ} (takenoko) — a bamboo shoot}

\label{lesson:JA-C494-takenoko}



\begin{quote}
Before the new one: say the Japanese for a cucumber, then the Japanese for a mushroom.
\end{quote}

\subsection*{You'll want to know: \ja{たけのこ}}
\textbf{\ja{たけのこ}} — \emph{takenoko} — \textquotedblleft{}a bamboo shoot\textquotedblright{}.

\begin{quote}
\textbf{\ja{はるに} \ja{たけのこを} \ja{たべます。}} — \emph{haru ni takenoko o tabemasu} — \textquotedblleft{}We eat bamboo shoots in spring.\textquotedblright{}
\end{quote}

\begin{cousinweb}
\textbf{\ja{たけのこ}} (\emph{takenoko}) is a native Japanese word (wago), inherited rather than borrowed from Chinese, and usually written in kana; its older origin is not settled (JMdict).
\end{cousinweb}

\begin{grammarlens}[title={What you've built}]
One more word for describing what you lived and felt.
\end{grammarlens}

\subsection*{Your turn}
\begin{itemize}
\raggedright
  \item \emph{Say it:} \emph{takenoko}
  \item \emph{Say it:} \emph{takenoko}, once more
  \item \emph{Say it:} \emph{haru ni takenoko o tabemasu} — We eat bamboo shoots in spring.
  \item \emph{Recall:} say the Japanese for a cucumber, then the Japanese for a mushroom, then say \emph{takenoko} again
\end{itemize}

\begin{tcolorbox}[breakable,colback=teal!4,colframe=teal!35!black,title={Before you move on}]
What does \ja{たけのこ} mean? (\textquotedblleft{}A bamboo shoot\textquotedblright{}.) Say it once more.
\end{tcolorbox}
\section[\texorpdfstring{\ja{しいたけ} (shiitake)}{しいたけ (shiitake)}]{\ja{しいたけ} (shiitake) — a shiitake mushroom}

\label{lesson:JA-C494-shiitake}



\begin{quote}
Before the new one: say the Japanese for a mushroom, then the Japanese for a bamboo shoot.
\end{quote}

\subsection*{You'll want to know: \ja{しいたけ}}
\textbf{\ja{しいたけ}} — \emph{shiitake} — \textquotedblleft{}a shiitake mushroom\textquotedblright{}.

\begin{quote}
\textbf{\ja{しいたけを} \ja{みそしるに} \ja{いれます。}} — \emph{shiitake o misoshiru ni iremasu} — \textquotedblleft{}I put shiitake in the miso soup.\textquotedblright{}
\end{quote}

\begin{cousinweb}
\textbf{\ja{しいたけ}} (\emph{shiitake}) is a native Japanese word (wago), inherited rather than borrowed from Chinese, and usually written in kana; its older origin is not settled (JMdict).
\end{cousinweb}

\begin{grammarlens}[title={What you've built}]
One more word for describing what you lived and felt.
\end{grammarlens}

\subsection*{Your turn}
\begin{itemize}
\raggedright
  \item \emph{Say it:} \emph{shiitake}
  \item \emph{Say it:} \emph{shiitake}, once more
  \item \emph{Say it:} \emph{shiitake o misoshiru ni iremasu} — I put shiitake in the miso soup.
  \item \emph{Recall:} say the Japanese for a mushroom, then the Japanese for a bamboo shoot, then say \emph{shiitake} again
\end{itemize}

\begin{tcolorbox}[breakable,colback=teal!4,colframe=teal!35!black,title={Before you move on}]
What does \ja{しいたけ} mean? (\textquotedblleft{}A shiitake mushroom\textquotedblright{}.) Say it once more.
\end{tcolorbox}
